% FRAGMENT A02 (drafter A: class model and main results)
% REPLACES: main.tex lines 158–177 ("The mechanism runs through two claims" to "… worth the cost of trading.")
% CHANGE: the two-value explanation becomes the class explanation (the split of a fixed value gap).
%   Lines 173–177 also contain investor sentences. They are reproduced verbatim except
%   "knows the challenger's value" -> "knows the challenger's class". If the institutional
%   drafter also edits these investor sentences, keep that drafter's investor wording and
%   keep the word "class".
% LABELS: none kept, removed, or added.
% NUM PLACEHOLDERS: none.
The mechanism runs through two claims on the same acquisition surplus.
The challenger holds the acquisition claim, what it keeps if it
prepares, bids, and wins. Target shareholders hold the traded claim,
what the auction pays them. The challenger's value comes from a
high-value or a low-value quality class. Each bidder bids its value in
the cash second-price auction. A prepared challenger's value then goes
to the two claims: a win splits it into the payment and the profit, and
a loss makes the winning incumbent pay the challenger's bid. The
expected value gap between the classes is fixed by the classes alone.
A stronger incumbent changes only who receives it. It outbids the
challenger more often and pays the challenger's bid, so more of the gap
reaches target proceeds. A winning challenger pays more, so less of
the gap remains as acquisition profit, at every belief. Target shares become
more sensitive to information about the challenger at the same time as
the challenger's own return to preparation falls. An investor who knows
the challenger's class trades on that sensitivity against noise demand.
Competitive market makers price order flow and anticipate the
preparation that each price will induce, but noise leaves the investor a
residual advantage. A larger gap in target proceeds can make that
advantage worth the cost of trading.
